\chapter{Itching in Pregnancy and Intrahepatic Cholestasis}

\section*{Definition}
Itching in pregnancy is common from skin stretching and dryness, but intense itching without a rash, especially on the palms and soles and worse at night, may indicate intrahepatic cholestasis of pregnancy.

\section*{When to See a Doctor}
Contact the maternity team promptly for severe unexplained itching, dark urine, pale stools, jaundice, right-upper abdominal pain, or reduced fetal movement. Arrange prompt evaluation for persistent, recurrent, unexplained, or function-limiting symptoms.

\section*{How It Develops}
The symptom results from irritation, inflammation, altered organ function, injury, infection, or disruption of normal neurologic or vascular signaling. Age, pregnancy status, onset, distribution, severity, and associated findings determine the urgency and likely causes.

\section*{Causes}
\begin{itemize}
  \item Dry skin, eczema, allergic reactions, scabies, obstetric cholestasis, liver disease, and medication reactions.
  \item Medication effects, environmental exposures, and underlying chronic disease may also contribute.
  \item Serious causes are less common but must be considered when symptoms are sudden, progressive, severe, or associated with systemic illness.
\end{itemize}

\section*{Related Symptoms}
\begin{itemize}
  \item Pain, fever, fatigue, nausea, weakness, or changes in normal function
  \item Bleeding, swelling, discharge, breathing difficulty, or neurologic symptoms when a serious cause is present
\end{itemize}

\section*{Medical Evaluation}
\begin{itemize}
  \item History addresses onset, duration, severity, triggers, injuries, exposures, medications, medical conditions, age, and pregnancy status when relevant.
  \item Examination includes vital signs and focused assessment of the relevant organ system, neurologic status, hydration, and function.
  \item Testing may include blood or urine studies, cultures, imaging, physiologic testing, fetal or pediatric assessment, or specialist examination.
\end{itemize}

\section*{Treatment Options}
Treatment is directed at the cause and may include supportive care, targeted medication, treatment of infection or inflammation, correction of metabolic disease, physical therapy, or specialist procedures. Persistent symptoms require reassessment.

\section*{Self-Care and Home Management}
Avoid known triggers, protect the affected area, maintain appropriate hydration, and record timing and associated symptoms. Use only age- or pregnancy-appropriate medicines recommended by a clinician. Do not delay emergency care while trying home treatment.

\section*{References}
\begin{thebibliography}{99}
\bibitem{itching_in_pregnancy:rcog} Royal College of Obstetricians and Gynaecologists. Intrahepatic cholestasis of pregnancy. RCOG Green-top Guideline No. 43; 2011 (updated 2022).
\bibitem{itching_in_pregnancy:bmj} Geenes V, Chappell LC. Intrahepatic cholestasis of pregnancy. \textit{BMJ}. 2014;349:g4712.
\end{thebibliography}
